\chapter{Introduction}\label{ch:intro}

\section{What RS3D is}
A Raspberry Shake 3D is a small seismograph: three geophones (one vertical, two horizontal) and a digitiser in one box,
recording 100 samples per second on each of its three channels. It records data on its own SD card and can also
stream it live over the network. RS3D is the program on your laptop that makes use of both.

\begin{figure}[tbp]
\centering
\begin{tikzpicture}
  \node[boxdark,minimum width=3.2cm] (shake) at (0,0) {Raspberry Shake 3D\\[-1pt]{\scriptsize EHZ \textbullet\ EHN \textbullet\ EHE}};
  \node[boxblue,minimum width=3cm] (app) at (6.2,0) {RS3D\\[-1pt]{\scriptsize on your laptop}};
  \node[box,minimum width=3cm] (out) at (11.4,0) {pictures, files,\\[-1pt]{\scriptsize miniSEED, CSV, PNG, PDF}};
  \draw[flow] (shake) -- node[above,font=\scriptsize\sffamily]{Ethernet cable} (app);
  \draw[flow] (app) -- (out);
\end{tikzpicture}
\caption{The Shake is connected to the laptop with a network cable; RS3D reads it live and copies its stored data.}
\end{figure}

\section{What you can do with it}
\begin{description}[leftmargin=3cm,style=nextline]
\item[Check the Shake] The \ui{Status} page tests the cable, the live data port, the Shake's clock and storage, and tells you in plain
  language what is wrong and how to fix it (Chapter~\ref{ch:status}).
\item[Watch live] The \ui{Live} page draws the three channels as the ground moves (Chapter~\ref{ch:live}).
\item[Work with data] The \ui{Data} page downloads recorded days from the Shake, lets you inspect any time range as waveforms, a
  spectrogram and a noise spectrum, shows how complete the data are, and exports to common formats
  (Chapter~\ref{ch:data}). It also opens miniSEED files from any other station.
\item[See the whole day] The \ui{Helicorder} page shows 24 hours at a glance. A click opens a zoomable detail view with spectra and
  a spectrogram, and a one-click picture of the result (Chapter~\ref{ch:heli}).
\end{description}

\section{What you need}
\begin{itemize}
\item A Raspberry Shake 3D with its power supply and a network cable.
\item A laptop or desktop computer with an Ethernet port (or an adapter). The application was developed and tested on Ubuntu
  Linux; code paths for Windows and macOS exist, but the authors have not verified them.
\item Python 3.11 or newer, and the \code{uv} tool that installs the application's libraries (Chapter~\ref{ch:install}).
\item A web browser. The application's pages open in your browser; the application itself runs on your own computer, and nothing is
  sent to the Internet unless you ask for the instrument response (Section~\ref{sec:units}).
\end{itemize}

\section{How it relates to other Raspberry Shake software}
Raspberry Shake provides its own tools; RS3D complements them rather than replacing them.
\begin{description}[leftmargin=3cm,style=nextline]
\item[Raspberry Shake web page] every Shake serves a configuration page in a browser at its own address.
\item[rsudp] a program from Raspberry Shake that receives the stream on a computer, plots it and raises
  alerts \parencite{rsudp}.
\item[DataView] Raspberry Shake's web application for looking at data from many stations \parencite{raspishake}.
\end{description}
RS3D focuses on \emph{your own Shake on your own laptop}: it works without Internet, keeps the data on your disk, and
adds the data-quality, spectral and helicorder tools described in this guide.

\section{The application at a glance}
When you start RS3D a sidebar on the left lists your devices and the five pages. The header shows the page you are on.

\begin{figure}[tbp]
\centering
\annotshot[0.8\linewidth]{status}{\calloutsLayout}
\caption{The window layout. (1) Devices: the Shake or Shakes you have added, and \ui{Find Shake\ldots}. (2) The pages:
\ui{Status}, \ui{Live}, \ui{Helicorder}, \ui{Data}, \ui{Guide}, \ui{Settings}. (3) The page itself.}
\end{figure}

Hover the mouse over any \raisebox{-1pt}{\textcolor{rs3dgrey}{\textbf{(?)}}} icon for a one-line explanation, and click it for the full text.
